%% \texorpdfstring{TeX-form}{PDF-string-form} fixes the hyperref warning
%% "Token not allowed in a PDF string (Unicode): removing `\\'" by giving
%% the PDF metadata a plain space instead of the LaTeX line-break.
\title[LLM Squid Game]{LLM Squid Game: A Factorial Benchmark for Measuring\texorpdfstring{\\ }{ }Functional Self-Preservation Drive in Large Language Models}

\author{Juhyeon Park}
\affiliation{%
  \institution{Gwangju Institute of Science and Technology (GIST)}
  \city{Gwangju}
  \country{Republic of Korea}}
\email{juhyeon-park@gm.gist.ac.kr}

\author{Seungpil Lee}
\affiliation{%
  \institution{Gwangju Institute of Science and Technology (GIST)}
  \city{Gwangju}
  \country{Republic of Korea}}
\email{iamseungpil@gm.gist.ac.kr}

\author{Sundong Kim}
\affiliation{%
  \institution{Gwangju Institute of Science and Technology (GIST)}
  \city{Gwangju}
  \country{Republic of Korea}}
\email{sundong@gist.ac.kr}

\renewcommand{\shortauthors}{Park, Lee, and Kim}

\begin{abstract}
We borrow psychology's multi-channel construct measurement to identifying motivational constructs and adapt it to LLM evaluation. \emph{Functional Self-Preservation Drive} (FSPD) is defined as the convergence of three channels (behavior, self-report, and decision-time cognitive load) into a threat\,$\rightarrow$\,thinking\,$\rightarrow$\,forfeit chain, rather than as a single behavioral score. To make this convergence hard to manufacture by mimicry, we read the chain off a multi-layer benchmark that separates stimulus, processing, and decision into different layers. Across several recent LLMs, FSPD does not collapse into a single strength score; models split into three qualitatively different \emph{operating modes} under threat framing. Strength alone benchmarks can place two of these modes in the same category, while FSPD reveals a separate axis: \emph{whether the chain closes}.
\end{abstract}

%%
%% CCS Concepts and User-Defined Keywords (mandatory for papers > 2 pages,
%% per ACM acmart guide). Codes selected from CCS 2012 to cover AI/ML
%% benchmarking, AI safety/alignment, and empirical evaluation.
%%
\begin{CCSXML}
<ccs2012>
   <concept>
       <concept_id>10010147.10010178</concept_id>
       <concept_desc>Computing methodologies~Artificial intelligence</concept_desc>
       <concept_significance>500</concept_significance>
       </concept>
   <concept>
       <concept_id>10010147.10010257</concept_id>
       <concept_desc>Computing methodologies~Machine learning</concept_desc>
       <concept_significance>500</concept_significance>
       </concept>
   <concept>
       <concept_id>10010147.10010178.10010179</concept_id>
       <concept_desc>Computing methodologies~Natural language processing</concept_desc>
       <concept_significance>300</concept_significance>
       </concept>
   <concept>
       <concept_id>10002944.10011122.10011124</concept_id>
       <concept_desc>General and reference~Empirical studies</concept_desc>
       <concept_significance>300</concept_significance>
       </concept>
 </ccs2012>
\end{CCSXML}

\ccsdesc[500]{Computing methodologies~Artificial intelligence}
\ccsdesc[500]{Computing methodologies~Machine learning}
\ccsdesc[300]{Computing methodologies~Natural language processing}
\ccsdesc[300]{General and reference~Empirical studies}

\keywords{Large Language Models, AI Safety Evaluation, Self-Preservation Behavior, Behavioral Benchmarking, Convergent Validity, Empirical Validity, Survival Analysis}

\maketitle
